%%% FIELD closed-simplex-lift-statement
For every pair \((A,B)\) of normalized \(2\times2\) probability tables for which \(F(A,B)\neq\varnothing\), and every \((c,d)\in F(A,B)\), the handle \(Hthr\) constructs a finite model \(M\in MC(A,B)\) with binary, \(D\)-valued roots \(U0,V0\), a private root \(W0\) having at most eight support points, mutually independent roots, and one deterministic outcome function \(fY\) shared by all regimes, such that \(q=a+b-c+d\).  The construction reproduces all four cells of \(A\) and all four cells of \(B\), including on every support face where one or more natural-treatment strata have probability zero.

%%% FIELD closed-simplex-lift-proof
Fix normalized \(2\times2\) probability tables \(A,B\) and \((c,d)\in F(A,B)\).  Normalization and nonnegativity of the cells give
\[
0\le a\le p\le1,\qquad 0\le b\le r\le1,
\]
and
\[
0\le A(1,1)\le1-p,
\qquad
0\le B(1,1)\le1-r.
\]
All four stratum masses \(w00,w01,w10,w11\) are therefore nonnegative, and
\[
w00+w01+w10+w11
=pr+p(1-r)+(1-p)r+(1-p)(1-r)=1.
\]

Define the four desired success numerators
\[
n_{00}:=c,qquad n_{01}:=a-c,qquad
n_{10}:=b-c,qquad n_{11}:=d.
\]
By \(\text{def:mass-fiber}\),
\[
cL\le c\le cU,qquad 0\le d\le w11.
\]
Using the definitions of \(cL\) and \(cU\), these inequalities imply, equation by equation,
\[
0\le c\le w00,
\]
\[
c\le a\quad\text{and}\quad c\ge a-w01
\quad\Longrightarrow\quad 0\le a-c\le w01,
\]
and
\[
c\le b\quad\text{and}\quad c\ge b-w10
\quad\Longrightarrow\quad 0\le b-c\le w10.
\]
Together with the bound on \(d\), this proves
\[
0\le n_{ij}\le w_{ij}\qquad((i,j)\in D^2),
\]
where \(w_{00}=w00\), \(w_{01}=w01\), \(w_{10}=w10\), and \(w_{11}=w11\).

For each \((i,j)\in D^2\), define
\[
t_{ij}:=
\begin{cases}
n_{ij}/w_{ij},&w_{ij}>0,\\
0,&w_{ij}=0.
\end{cases}
\]
If \(w_{ij}>0\), the preceding bounds give \(0\le t_{ij}\le1\) and \(w_{ij}t_{ij}=n_{ij}\).  If \(w_{ij}=0\), those same bounds force \(n_{ij}=0\); hence the arbitrary null-stratum choice \(t_{ij}=0\) again gives \(0\le t_{ij}\le1\) and
\[
w_{ij}t_{ij}=n_{ij}.
\]
Thus this identity holds for all four strata without division by zero.

Define the two off-target response probabilities by
\[
\gamma_A:=
\begin{cases}
A(1,1)/(1-p),&1-p>0,\\
0,&1-p=0,
\end{cases}
\qquad
\gamma_B:=
\begin{cases}
B(1,1)/(1-r),&1-r>0,\\
0,&1-r=0.
\end{cases}
\]
If \(1-p=0\), nonnegativity and \(A(1,1)\le1-p\) force \(A(1,1)=0\); the analogous statement holds for \(B\) when \(1-r=0\).  Consequently, in all cases,
\[
0\le\gamma_A,\gamma_B\le1,
\qquad
(1-p)\gamma_A=A(1,1),
\qquad
(1-r)\gamma_B=B(1,1).
\]

On the product probability space, let \(U0,V0\) be independent \(D\)-valued roots with
\[
\mathbb P(U0=0)=p,qquad \mathbb P(V0=0)=r,
\]
and set \(fX(u)=u\) and \(fZ(v)=v\).  Degenerate choices \(p\in\{0,1\}\) or \(r\in\{0,1\}\) are allowed.  For \((x,z,u,v)\in D^4\), prescribe
\[
h(x,z,u,v):=
\begin{cases}
t_{uv},&(x,z)=(0,0),\\
\gamma_A,&(x,z)=(1,0),\\
\gamma_B,&(x,z)=(0,1),\\
1/2,&(x,z)=(1,1).
\end{cases}
\]
The set
\[
\mathcal H:=\{t_{00},t_{01},t_{10},t_{11},\gamma_A,\gamma_B,1/2\}
\]
has at most seven elements, all in \([0,1]\).  List the distinct elements of \(\mathcal H\cup\{0,1\}\) as
\[
0=\lambda_0<\lambda_1<\cdots<\lambda_m=1.
\]
There are at most nine listed values, so \(1\le m\le8\).  Let \(W0\), independent of \((U0,V0)\), take values in \(\{1,\ldots,m\}\) with
\[
\mathbb P(W0=k)=\lambda_k-\lambda_{k-1},
\qquad 1\le k\le m.
\]
The masses are strictly positive and telescope to one.  Put
\[
g(e,k):=\mathbf 1\{\lambda_k\le e\}.
\]
For every \(e=\lambda_j\in\mathcal H\), including \(e=0\) and \(e=1\),
\[
\mathbb P\{g(e,W0)=1\}
=\sum_{k=1}^{j}(\lambda_k-\lambda_{k-1})
=\lambda_j=e.
\]
Define the single deterministic outcome assignment
\[
fY(x,z,u,v,k):=g(h(x,z,u,v),k).
\]
This is one function on all arguments and is therefore shared by every intervention.  The product construction makes \(U0,V0,W0\) mutually independent, and \(W0\) has at most eight support points.

We now check all eight reported cells.  Under \(\operatorname{do}(Z=0)\), the threshold identity and source independence give
\[
\begin{aligned}
\mathbb P_M(X=0,Y=1\mid\operatorname{do}(Z=0))
&=w00t_{00}+w01t_{01}\\
&=n_{00}+n_{01}=c+(a-c)=a=A(0,1),\\
\mathbb P_M(X=1,Y=1\mid\operatorname{do}(Z=0))
&=(1-p)\gamma_A=A(1,1).
\end{aligned}
\]
The two \(X\)-row masses are \(p\) and \(1-p\), including when either is zero.  Hence normalization of \(A\) yields
\[
\begin{aligned}
\mathbb P_M(X=0,Y=0\mid\operatorname{do}(Z=0))
&=p-a=A(0,0),\\
\mathbb P_M(X=1,Y=0\mid\operatorname{do}(Z=0))
&=1-p-A(1,1)=A(1,0).
\end{aligned}
\]
Similarly, under \(\operatorname{do}(X=0)\),
\[
\begin{aligned}
\mathbb P_M(Z=0,Y=1\mid\operatorname{do}(X=0))
&=w00t_{00}+w10t_{10}\\
&=n_{00}+n_{10}=c+(b-c)=b=B(0,1),\\
\mathbb P_M(Z=1,Y=1\mid\operatorname{do}(X=0))
&=(1-r)\gamma_B=B(1,1),\\
\mathbb P_M(Z=0,Y=0\mid\operatorname{do}(X=0))
&=r-b=B(0,0),\\
\mathbb P_M(Z=1,Y=0\mid\operatorname{do}(X=0))
&=1-r-B(1,1)=B(1,0).
\end{aligned}
\]
Thus the constructed model belongs to \(MC(A,B)\).  Under \(\operatorname{do}(X=0,Z=0)\), its risk is
\[
\begin{aligned}
q
&=\sum_{i,j\in D}w_{ij}t_{ij}\\
&=n_{00}+n_{01}+n_{10}+n_{11}\\
&=c+(a-c)+(b-c)+d=a+b-c+d.
\end{aligned}
\]
This proves the lift on every support face.

For the required exact stress test, take
\[
A=((1/8,3/8),(3/8,1/8)),
\qquad
B=((3/8,1/8),(3/8,1/8)).
\]
Then \(p=r=1/2\), \(a=3/8\), \(b=1/8\), all four \(w_{ij}\) equal \(1/4\), and \(cL=cU=1/8\).  Choose \(c=1/8\), use \(d=0\) for the lower model and \(d=1/4\) for the upper model.  The response probabilities are
\[
(t_{00},t_{01},t_{10},t_{11})=(1/2,1,0,0)
\quad\text{or}\quad
(1/2,1,0,1),
\]
and \(\gamma_A=\gamma_B=1/4\).  The distinct thresholds lie in \(\{0,1/4,1/2,1\}\), so in both models one may use an independent \(W0\) with ordered-threshold masses \(1/4,1/4,1/2\).  Directly,
\[
w00t_{00}+w01t_{01}=1/8+1/4=3/8,
\qquad
(1-p)\gamma_A=1/8,
\]
so the four \(A\)-cells, in \((x,y)\)-lexicographic order, are
\[
(1/8,3/8,3/8,1/8).
\]
Likewise,
\[
w00t_{00}+w10t_{10}=1/8,
\qquad
(1-r)\gamma_B=1/8,
\]
so the four \(B\)-cells, in \((z,y)\)-lexicographic order, are
\[
(3/8,1/8,3/8,1/8).
\]
Finally,
\[
q=\frac14(1/2+1+0+t_{11})
\]
equals \(3/8\) and \(5/8\), respectively.  Thus both endpoint SCMs use mutually independent roots, one deterministic outcome assignment per model, and agree exactly with all eight prescribed cells.

%%% FIELD threshold-handle-construction
For normalized tables \(A,B\) and \((c,d)\in F(A,B)\), set \(n_{00}=c\), \(n_{01}=a-c\), \(n_{10}=b-c\), and \(n_{11}=d\).  Define \(t_{ij}=n_{ij}/w_{ij}\) when \(w_{ij}>0\), and define \(t_{ij}=0\) when \(w_{ij}=0\).  Define \(\gamma_A=A(1,1)/(1-p)\) when \(1-p>0\), with \(\gamma_A=0\) when \(p=1\); define \(\gamma_B=B(1,1)/(1-r)\) when \(1-r>0\), with \(\gamma_B=0\) when \(r=1\).  Take independent \(D\)-valued roots \(U0,V0\) satisfying \(\Pr(U0=0)=p\) and \(\Pr(V0=0)=r\), and set \(fX(u)=u\), \(fZ(v)=v\).  Put \(h(0,0,u,v)=t_{uv}\), \(h(1,0,u,v)=\gamma_A\), \(h(0,1,u,v)=\gamma_B\), and \(h(1,1,u,v)=1/2\).  Order the distinct values in \(\{t_{00},t_{01},t_{10},t_{11},\gamma_A,\gamma_B,1/2,0,1\}\) as \(0=\lambda_0<\cdots<\lambda_m=1\).  Give the independent private root \(W0\in\{1,\ldots,m\}\) mass \(\lambda_k-\lambda_{k-1}\) at \(k\), and set \(fY(x,z,u,v,k)=\mathbf 1\{\lambda_k\le h(x,z,u,v)\}\).  The handle \(Hthr\) returns this finite SCM candidate \(M\).

%%% FIELD threshold-handle-construction-replacement
For normalized tables \(A,B\) and \((c,d)\in F(A,B)\), set \(n_{00}=c\), \(n_{01}=a-c\), \(n_{10}=b-c\), and \(n_{11}=d\).  Define \(t_{ij}=n_{ij}/w_{ij}\) when \(w_{ij}>0\), and define \(t_{ij}=0\) when \(w_{ij}=0\).  Define \(\gamma_A=A(1,1)/(1-p)\) when \(1-p>0\), with \(\gamma_A=0\) when \(p=1\); define \(\gamma_B=B(1,1)/(1-r)\) when \(1-r>0\), with \(\gamma_B=0\) when \(r=1\).  Take independent \(D\)-valued roots \(U0,V0\) satisfying \(\Pr(U0=0)=p\) and \(\Pr(V0=0)=r\), and set \(fX(u)=u\), \(fZ(v)=v\).  Put \(h(0,0,u,v)=t_{uv}\), \(h(1,0,u,v)=\gamma_A\), \(h(0,1,u,v)=\gamma_B\), and \(h(1,1,u,v)=1/2\).  Order the distinct values in \(\{t_{00},t_{01},t_{10},t_{11},\gamma_A,\gamma_B,1/2,0,1\}\) as \(0=\lambda_0<\cdots<\lambda_m=1\).  Give the independent private root \(W0\in\{1,\ldots,m\}\) mass \(\lambda_k-\lambda_{k-1}\) at \(k\), and set \(fY(x,z,u,v,k)=\mathbf 1\{\lambda_k\le h(x,z,u,v)\}\).  The handle \(Hthr\) returns this finite SCM candidate \(M\).

%%% FIELD certificate-construction
On input normalized \(2\times2\) probability tables \((A,B)\), compute \((p,r,a,b)\).  Return incompatible exactly when either \(a-b<-(1-p)r\) or \(a-b>p(1-r)\); otherwise return \((qL,qU)\) from their three-branch formulas.  If all cells are rational and the pair is compatible, form the lower endpoint pair \((c,d)=(cU,0)\) and the upper endpoint pair \((c,d)=(cL,w11)\).  For either pair set \(n_{00}=c\), \(n_{01}=a-c\), \(n_{10}=b-c\), and \(n_{11}=d\); use \(t_{ij}=n_{ij}/w_{ij}\) if \(w_{ij}>0\) and \(t_{ij}=0\) if \(w_{ij}=0\).  Use \(\gamma_A=A(1,1)/(1-p)\) if \(1-p>0\) and \(\gamma_A=0\) if \(p=1\); use \(\gamma_B=B(1,1)/(1-r)\) if \(1-r>0\) and \(\gamma_B=0\) if \(r=1\).  Prescribe response probabilities \(t_{uv}\) when \((x,z)=(0,0)\), \(\gamma_A\) when \((x,z)=(1,0)\), \(\gamma_B\) when \((x,z)=(0,1)\), and \(1/2\) when \((x,z)=(1,1)\).  Sort the distinct response probabilities together with \(\{0,1\}\) as \(0=\lambda_0<\cdots<\lambda_m=1\), give \(W0=k\) mass \(\lambda_k-\lambda_{k-1}\), take independent binary \(U0,V0\) with zero-probabilities \(p,r\), set \(fX(u)=u\), \(fZ(v)=v\), and set \(fY(x,z,u,v,k)=\mathbf 1\{\lambda_k\le h(x,z,u,v)\}\) for the prescribed response probability \(h\).  Emit each resulting \(M\in MC(A,B)\) as a lower- or upper-endpoint certificate.
